\subsection{Equações Diferenciais da Cinemática}
\label{sec:equacoes_cinematica}

A atitude da espaçonave visitante foi descrita utilizando os quaternions unitários por conta da ausência de singularidades (\emph{gimbal lock}) e maior eficiência computacional em simulações numéricas.

\subsubsection{Quaternions unitários}

Um quaternion unitário é definido como:

\begin{equation}
\label{eq:quaternion_unitario_definicao}
\mathbf{q} =
\begin{bmatrix}
q_0 \\ q_1 \\ q_2 \\ q_3
\end{bmatrix}
=
\begin{bmatrix}
\cos\frac{\alpha}{2} \\
u_x \sin\frac{\alpha}{2} \\
u_y \sin\frac{\alpha}{2} \\
u_z \sin\frac{\alpha}{2}
\end{bmatrix},
\end{equation}

onde $\alpha$ é o ângulo de rotação e $\vec{u}=(u_x,u_y,u_z)^T$ é o eixo unitário.
A condição de normalização é:

\begin{equation}
\label{eq:normalizacao_quaternion}
\lVert \mathbf{q} \rVert^2 = q_0^2+q_1^2+q_2^2+q_3^2 = 1.
\end{equation}

\subsubsection{Cinemática diferencial}

Considerando a velocidade angular $\vec{\omega}$ expressa no corpo da espaçonave, a cinemática do quaternion é dada por:

\begin{equation}
\label{eq:qdot}
\dot{\vec{q}} = \tfrac{1}{2}\,\Omega(\vec{q})\,\vec{\omega},
\end{equation}

sendo:

\begin{equation}
\Omega(\vec{q}) =
\begin{bmatrix}
-{\vec \epsilon\,}^T \\
q_0 \mathbf{I}_3 + [\vec \epsilon\,]_\times
\end{bmatrix},
\qquad
\vec{q} = \begin{bmatrix} q_0 \\
\vec{\epsilon} \end{bmatrix},
\end{equation}
onde $[\vec \epsilon\,]_\times$ é a matriz antissimétrica associada ao vetor $\vec \epsilon$.

\subsubsection{Acoplamento com a dinâmica}

No caso de corpo rígido, a equação \eqref{eq:qdot} descreve a evolução temporal da atitude a partir da velocidade angular. 
No entanto, quando o MR está em movimento.
Como visto em \eqref{eq:dinamica_acoplada}, a velocidade angular $\vec{\omega}$ não é independente; ela resulta da integração das equações de Euler com inércia variável.
A solução das equações diferenciais da cinemática depende da evolução conjunta de $\vec{\omega}$, a qual é influenciada pelas variáveis generalizadas do MR.
Esse acoplamento implica que o modelo da atitude da espaçonave visitante não é o de um corpo rígido, mas de um sistema dinâmico composto, em que a matriz inversa de inércia fornece as equações diferenciais completas da cinemática.

A equação de Euler modificada, apresentada em \eqref{eq:dinamica_acoplada}, pode ser reescrita isolando $\dot{\vec{\omega}}$:

\begin{equation}
\label{eq:euler_modificada}
\dot{\vec{\omega}}
= \mathbf{I}^{-1}(q)\Big(\vec{M}_{controle} + \vec{M}_{manip} - \vec{\omega} \times \big(\mathbf{I}(q)\vec{\omega}\big)\Big).
\end{equation}

onde $\mathbf{I}(q)$ é função dos ângulos e velocidades das juntas do manipulador robótico, de modo que tanto a matriz de inércia quanto sua inversa variam no tempo, sendo assim, a integração de \eqref{eq:qdot} e \eqref{eq:euler_modificada} fornece as equações diferenciais da cinemática da atitude da espaçonave visitante com a movimentação realizada pelo MR.
